% Actual perfect source factor and its canonical-centralizer lift.
\subsection{An independent perfect source factor in the full descent kernel}

Let $q_R:B\to R$ be any specified finite quotient.  Retain
the full six-coordinate kernel $K$ of B.2.13 and put
$K_R=K\cap\ker q_R$.  For a prime $p\geq5$, use the
actual modular quotient
\[
 q_p:B\twoheadrightarrow Q_p=\operatorname{PSL}_2(\mathbb F_p),
 \qquad s\longmapsto[S],\quad r\longmapsto[R_0],
 \quad
 S=\begin{pmatrix}0&-1\\1&0\end{pmatrix},\quad
 R_0=\begin{pmatrix}0&-1\\1&-1\end{pmatrix}.
 \tag{B.AB1}
\]
For every $p\nmid|R|$ define
\[
 U_R=B/K_R,\quad K_{R,p}=K_R\cap\ker q_p,\quad
 U_{R,p}=B/K_{R,p},\quad
 U_{R,p}^{\mathrm c}=B/\overline{[B,K_{R,p}]}.
 \tag{B.AB2}
\]
All maps are induced by the named source elements.

\begin{result}{B.2.29}{Proved}{A full positive lift with an independent perfect source frame}
For every pair $(R,p)$ specified above, the actual joint map is
\[
 U_{R,p}\xrightarrow{\ \sim\ }U_R\times Q_p.
 \tag{B.AB3}
\]
The pure element $b=(1,q_p(s))$ in this product has a
pure lift to $U_{R,p}^{\mathrm c}$ that centralizes the
named $s$.  There is consequently a full positive
cusp-fixed automorphism $\Phi$ of $B$ with pure frames
$\beta,A$ satisfying
\[
 \Phi_{U_{R,p}}=1,\qquad
 \beta_{U_{R,p}}=(1,q_p(s)),\qquad A_{U_{R,p}}=1.
 \tag{B.AB4}
\]
Its commutator normalization gives an actual
$a_{R,p}\in H\cap\ker D_P$.
\end{result}
\begin{proof}
The Brunnian word $w$ of (B.Z4) lies in the actual
discrete six-coordinate kernel $K_0$.  Under the integral
pure matrix representation $x\mapsto U^2$,
$y\mapsto L^2$, direct multiplication gives
\[
 w\longmapsto
 \begin{pmatrix}-191487&264064\\-721024&994305\end{pmatrix},
 \qquad
 \gcd(w-I)=128,\quad\gcd(w+I)=2.
 \tag{B.AB5}
\]
The two gcds mean the gcds of all four displayed
difference entries.  Thus the projective image of $w$
is nonidentity at every odd prime.  The integral matrix
and both gcds are retained in the exact checker
\path{check_prime_harmonic_cartans.py}.

The group $Q_p$ is nonabelian simple, hence perfect;
one primary proof is \cite[Theorems 3.3--3.5]{CA_ConradPSL}.
Surjectivity of (B.AB1) also follows directly because
the images contain the upper unipotents and their
$S$ conjugates, which generate $\operatorname{SL}_2(\mathbb F_p)$.
Consequently the normal subgroup $q_p(K_0)$ is all
of $Q_p$.  Put $K_{0,R}=K_0\cap\ker(q_R|_{B_0})$.
Its image under $q_p$ is again normal in $Q_p$.
A trivial image would make $Q_p$ a quotient of
$K_0/K_{0,R}$, a subgroup of the finite group $R$.
That would force $p$ to divide $|R|$, contrary to its
specified choice.  Thus
\[
                         q_p(K_{0,R})=Q_p.
 \tag{B.AB6}
\]
The joint image of $B_0$ in $U_{0,R}\times Q_p$ is
therefore the entire direct product, where
$U_{0,R}=B_0/K_{0,R}$.  The same argument with the
contained discrete subgroup proves (B.AB3) for the
actual profinite groups.

We prove the centralizer lift in the discrete source
first.  Set $K_{0,R,p}=K_{0,R}\cap\ker q_p$ and
$U_{0,R,p}=B_0/K_{0,R,p}=U_{0,R}\times Q_p$.
Since $K_{0,R}\subseteq B_0'$ and the quotient still
maps onto $U_0$, its abelianization is exactly $C_6$.
The direct product has the same abelianization because
$Q_p$ is perfect.  As $H_2(B_0,\mathbb Z)=0$, the
five-term homology sequence identifies its canonical
central kernel with
\[
 \begin{aligned}
 M_{0,R,p}
   &=K_{0,R,p}/[B_0,K_{0,R,p}]\\
   &\cong H_2(U_{0,R}\times Q_p,\mathbb Z)\\
   &\cong H_2(U_{0,R},\mathbb Z)\oplus H_2(Q_p,\mathbb Z).
 \end{aligned}
 \tag{B.AB7}
\]
The last isomorphism is induced by the two actual
projections and their coordinate sections.  In the
integral K\"unneth formula the only possible cross
term is $H_1(U_{0,R},\mathbb Z)\otimes H_1(Q_p,\mathbb Z)$,
which is zero.  No calculation of a Schur multiplier
or of its exponent is being inserted here.

Choose an actual word $v\in K_{0,R}$ with
$q_p(v)=q_p(s)$, using (B.AB6).  It represents
$b=(1,q_p(s))$, so $[v,s]$ belongs to the central
kernel in (B.AB7).  Its image there is trivial.
To verify this integrally, take any character of that
abelian kernel to $\mathbb Q/\mathbb Z$ and push out
the canonical central extension.  The universal-coefficient
sequence has no Ext term with these divisible
coefficients, so (B.AB7) says that the resulting
degree-two class is a sum of the pullbacks of classes
from the two factors.  In a class pulled back from
$U_{0,R}$, the element $b$ projects to one, so its
lift commutes with a lift of $s$.  In a class pulled
back from $Q_p$, the two elements $b,s$ project to
the same element $q_p(s)$.  Their lifts differ by a
central element and therefore commute as well.
The commutator evaluation is additive under the
sum of central-extension classes.  Every such
character thus kills the class of $[v,s]$.
Characters to $\mathbb Q/\mathbb Z$ separate elements
of an abelian group: a character detecting a
nonzero element of its cyclic subgroup extends
because $\mathbb Q/\mathbb Z$ is injective.
It follows that
\[
                         [v,s]\in[B_0,K_{0,R,p}].
 \tag{B.AB8}
\]

This word identity maps directly into
$\overline{[B,K_{R,p}]}$.  Therefore the image
$\widetilde b$ of $v$ in the full central cover
of (B.AB2) is a pure element centralizing $s$.
Use the identity automorphism of that cover with
positive $s$-frame $\widetilde b$, positive
$r$-frame one, and cusp frame one.
The full central-cover image equality B.2.14,
including its prescribed frame values, gives
exactly (B.AB4).

Finally put $\beta^{\rm ab}=(u,v_0)$ and $t=v_0-u$.
Since $\beta\in K$, the first two target maps into
$B^{\rm ab}=C_6$ give $3u+v_0=u+3v_0=0$ modulo six.
Thus $t\in6\widehat{\mathbb Z}$, as in B.2.27.
The normalized automorphism
$\Phi_0=\operatorname{Inn}(x^t)\Phi$ gives
$a_{R,p}\in H$.  Before this normalization the
descended action is the identity, since
$\Phi_U=1$.  After normalization it changes by
$\operatorname{Inn}(\varphi_P(x)^t)
=\operatorname{Inn}(s^t)=1$.  Hence its literal
normalized descendant is the identity, proving
$a_{R,p}\in\ker D_P$.
\end{proof}

